\documentclass[openany]{book}
\usepackage[utf8]{inputenc}
\usepackage{amsmath}
\usepackage{amsthm}
\usepackage{array}
\usepackage{biblatex}
\usepackage{latexsym,amsmath,xcolor,multicol,booktabs,calligra}

\usepackage{enumitem}
\usepackage{tabularx}
\usepackage{amssymb}
\usepackage{mathtools}
\usepackage{mathrsfs}
\usepackage{tikz}
\usepackage[margin=1in,bindingoffsest=0cm]{geometry}
\usepackage{circuitikz}
\usetikzlibrary{positioning}
\usepackage{caption}
\usepackage{xcite}
\usepackage{fancyhdr}
\usepackage{relsize}
\usepackage{lmodern}
\usepackage{slantsc}
\usepackage{hyperref}
\usepackage{cancel}
\usepackage{subcaption}
\usepackage{nicematrix}
\usetikzlibrary{shapes}
\usepackage{float}
\usepackage{quiver}

% Theorems-----------------------------------------------
\newtheorem{thm}{Theorem}[section]
\newtheorem{lem}[thm]{Lemma}
\newtheorem{cor}[thm]{Corollary}
\newtheorem{prop}[thm]{Proposition}
\newtheorem{conj}[thm]{Conjecture}
\newtheorem{answer}[thm]{Answer}
\theoremstyle{definition}
\newtheorem{prob}[thm]{Problem}
\newtheorem{obs}[thm]{Observation}
\newtheorem{definition}[thm]{Definition}
\newtheorem{rem}[thm]{Remark}
\newtheorem{question}[thm]{Question}
\newtheorem{nota}[thm]{Notation}


\newtheorem{example}[thm]{Example}
\AtBeginEnvironment{example}{\pushQED{\qed}\renewcommand{\qedsymbol}{$\triangle$}}
\AtEndEnvironment{example}{\popQED\endexample}

\newtheorem{procedure}[thm]{Procedure}
\AtBeginEnvironment{procedure}{ \pushQED{\qed}\renewcommand{\qedsymbol}{$\triangle$}}
\AtEndEnvironment{procedure}{\popQED\endprocedure}

\pagestyle{fancy}
\fancyhf{}
\fancyhead[RO]{\nouppercase{\rightmark}}
\fancyhead[LE]{\nouppercase{\leftmark}}
\fancyhead[RE,LO]{\thepage}

\title{Lecture Notes on Combinatorics: Draft}
\author{Garrett J. Kepler}
\date{Updated 2.4.26}
\makeindex

\begin{document}
\maketitle


\chapter{What is combinatorics?}

One great property of the field of combinatorics is most people have already encountered combinatorial questions in their own life.
\begin{center}
\emph{Considering the optimal path we should take to get groceries and visit the coffee shop then return home. Now that we have the groceries, figuring out which set of dishes we can cook for the party when we've invited a vegetarian, a non-vegetarian, and someone allergic to peanuts. Now that we know what to cook, determining the best drinks paired with which dishes.}
\end{center}
As simple of an example party planning is, we have described three combinatorial problems already: minimal path, dish enumeration, and drink/dish combinations. \\
Many of the questions we ask in combinatorics focus on counting of some sort. 


\begin{itemize}
    \item Enumeration: How many ways$\dots$?
    \item Existence: Does there exist$\dots$?
    \item Extremal: What is the largest/smallest possible$\dots$?
    \item Expectation: What is the expected number of$\dots$?
\end{itemize}

\subsection{Week's Warmups}

\begin{question}\hfill

    How many ways are there to assign 3 students 3 distinct Math 325 projects? 4 students 4 distinct projects? $n$ students $n$ distinct projects?
\end{question}

\begin{question}\hfill

    How many people must attend the first day of lecture in Math 325 to guarantee 2 people either both have met before or both have not met before?
\end{question}

\begin{question}\hfill

    How many people must attend the first day of lecture in Math 325 to guarantee 3 people either have all met before or have all not met before?
\end{question}
These preliminary warmup problems are classic combinatorial questions. The first, a question of enumeration, we may have a straightforward way to answer already:
\begin{answer}\hfill

    Well, we could assign student A project 1. Or, we could assign them project 2. Or, we could assign them project 3. That's 3 ways. Okay, for student B, we could assign project 1. Or, we could assign them project 2. Or,$\dots$ This is exhausting listing them out. I'm bored.
\end{answer}
What we are going to learn in this class is methods to avoid this one by one counting. Instead, we will ``count''. Rather than assigning a specific student to a specific project one by one, we'll deal with the collection of students and the collection of projects in their entirety:
\begin{answer}\hfill

    Given a student, we assign them 1 of the 3 projects. Since we want distinct assignments, the next student must be assigned one of the other 2 projects. Likewise, the final student must be assigned the leftover project. Using what we will later refer to as the \emph{multiplication principle}, we have $3\cdot 2\cdot 1=6$ ways to do this.
\end{answer}
Questions 1.2 and 1.3 ask questions of guaranteed substructures that we will discuss later in the Ramsey Theory section.

\chapter{Problem Solving Basics \& Set Theory}
We will start with questions of \emph{enumeration} and dedicate a fair amount of time to this concept.

\begin{definition}[Enumeration]\hfill

    The problem of \emph{enumeration} is to count the total number of objects with a specified property.
\end{definition}
\begin{example}\hfill

    A few basic examples:
    \begin{itemize}
        \item How many ways are there to assign $n$ students to $m$ projects?
        \item How many distinct Sudoku puzzles are there?
        \item How many total (potentially isomorphic) graphs are there are on $n$ nodes?
    \end{itemize}
\end{example}
This questions can often be difficult to answer given a certain object with a desired property. However, the philosophical approach we often take is: \textbf{simplify}!\\

Our approach to these problems can be summarized in a joke told to me by a faculty member at WSU:

\begin{center}
    \emph{A mathematician walks into a room. In one corner, he sees an empty bucket. In another corner, he sees a sink with a water faucet. In yet another corner, he sees a fire. Leaping into action, he picks up the bucket, fills it up at the faucet, and douses the fire.}\\
    
    \emph{The next day, the same mathematician returns to the room. Once again, he sees a fire in the same corner. But, this time the bucket rests next to the fire full of water. Once more he leaps into action: he picks up the bucket, dumps the water in the sink, places the bucket empty where it was from the day before and leaves. As he exits, he exclaims victoriously "I’ve reduced it to a previously solved problem!"}
\end{center}
Our first goal will be to mimic this proverbial mathematician's mentality: \emph{simplify}! Once we've simplified, we'll ask ourselves: ``can we generalize this?''\\

\subsection{Simplification}

\begin{question}\label{question1}\hfill

Consider the following problems:
    \begin{enumerate}
        \item How many whole numbers are between 1 and 325 inclusive?
        \item How many whole numbers are between 28 and 352 inclusive?   
        \item How many even numbers are there between 28 and 352 inclusive?
        \item How many numbers between 28 and 352 are divisible by 3?
    \end{enumerate}
\end{question}
On the one hand, these are four distinct problems. On the other, you may be able to see there is a process we can use to answer all of them. One main approach we often utilize in combinatorics is the
following: to count what we need to count, reframe the problem as something we already know how to count. To quote a great mathematician:
\begin{center}
    ``\emph{I have reduced it to a previously solved problem!}''
\end{center}
Consider our first problem: “How many whole numbers are between 1 and 325 inclusive?”
\begin{itemize}
    \item 1 is our $1^{st}$, 2 is our $2^{nd}$, and so on, 325 is our $325^{th}$. So, there must be 325 such numbers.
\end{itemize}
This was so easy to solve! What we want to do now is use this solved problem to answer another. Consider our second problem: “How many whole numbers are between 28 and 352 inclusive?”
\begin{itemize}
    \item We begin counting: 28 is our $1^{st}$ number, 29 is our $2^{nd}$, and so on, so 352 must be our $?^{th}$.
\end{itemize}
What do we put here? No idea. We only know how to count numbers between 1 and 325. Let's instead transform this counting problem to the one we know how to count:
\begin{itemize}
    \item We may notice 28 is 27 away from 1. Likewise, 29 is 27 away from 2 and so on.
    \item So, take each number $28,29,30,\dots, 352$ and subtract 27: $1,2,3,\dots, 325$.
    \item Amazing! We already know this answer: there are 325 such numbers. 
\end{itemize}
One issue you may already have thought of: ``how do I know there are the same amount of numbers between 1 and 325 as 28 and 352 by simply subtracting 27?'' That is a great point. We don't yet know. However, this is the beauty of a \emph{bijection}. Something we'll talk about: if we have a bijection between two sets, we know they are the same size. More on the formalities later!\\
\begin{center}
    "\emph{I have generalized the method to solve a previously unsolved
problem!}"
\end{center}
For now, consider the third problem: ``How many even numbers are there between 28 and 352 inclusive?''. Again, we don't know how to count this. But, on the second problem, we transformed to the numbers $1,2,3,\dots, 325$ and made our problem easier. Let's try something similar:
\begin{itemize}
    \item Take each number, $28,30,32,\dots,350, 352$, and divide by $2$: $14,15,16,\dots, 175, 176$.
    \item They may look different, but we can transform similarly to problem two. Subtract 13 and we get $1,2,3,\dots, 162, 163$.
    \item There are 163 such numbers, so there must be 163 such even numbers between 28 and 352.
\end{itemize}
Now that we have a specific method of solution, we might reformat this method to a more general problem. We may want to count the numbers between $m$ and $n$ ($m<n$). \\
In our second and third problem, we took each number and found a way to obtain $\{1,2,3,\dots, n\}$ so that we could enumerate the solutions. If we want to take our $1^{st}$ number $m$ and send it the first natural number $1$, we might subtract $m-1$ (since $m-(m-1)=m-m+1=1$). For each of these numbers, we get $m-(m-1),m+1-(m-1), m+2-(m-1),\dots, n-1-(m-1), n-(m-1)$ which really is $1,2,3,\dots, n-m, n-m+1$. As such, we obtain a count of $n-m+1$ integers between $m$ and $n$ ($m<n$). What's left to prove is that the process of subtracting $m-1$ here guarantees the two sets of numbers are the same size. The next section will discuss how to use a \emph{bijection} to do so.

\section{Set Theory Basics}
This concept of bijection we subtly utilized previously is actually a fundamental result combinatorialists love to exploit. Before we get there, some basics are necessary. 
\begin{definition}[Set]\hfill

    A \emph{set} is a collection of distinct objects. E.g. The set of integers between 28 and 56 = $\{28,29,30,\dots, 56\}$, a set of graphs = $\{ {
\includegraphics[height=2ex,valign=m]{math 593/photos/8.23.25/gseq1.png}},{\includegraphics[height=2ex,valign=m]{math 593/photos/8.23.25/gseq2.png}},{\includegraphics[height=2ex,valign=m]{math 593/photos/8.23.25/gseq3.png}},{\includegraphics[height=2ex,valign=m]{math 593/photos/8.23.25/gseq4.png}},\dots \}$, etc.
\end{definition}
Note in our definition of a set we require \emph{distinctness}. As such, the collection $\{1,1,2,3,\dots,29\}$ is not a set. Likewise with any collection with repeating elements, but this permits another important definition:
\begin{definition}[Multiset]\hfill

    A \emph{multiset} is a collection of not necessarily distinct objects. E.g. The multiset of integers = $\{28,28,28,29,30,\dots, 56\}$, the multiset of graphs = $\{ {
\includegraphics[height=2ex,valign=m]{math 593/photos/8.23.25/gseq1.png}},{
\includegraphics[height=2ex,valign=m]{math 593/photos/8.23.25/gseq1.png}},{\includegraphics[height=2ex,valign=m]{math 593/photos/8.23.25/gseq2.png}},{\includegraphics[height=2ex,valign=m]{math 593/photos/8.23.25/gseq3.png}},{\includegraphics[height=2ex,valign=m]{math 593/photos/8.23.25/gseq4.png}},\dots \}$, etc.
\end{definition}
However, we won't touch multisets until later on. For now, let us discuss relations between sets.
\begin{definition}[Function]\hfill

    A function $f:A\to B$ is a relation between sets $A$ and $B$ such that for each $a\in A$ ($a$ in $A$), there is a unique $b\in B$ such that $f(a)=b$.
\end{definition}

\begin{example}\hfill

As an example, on the left, we have $g:A\to B$ that assigns the second element in $A$ to two elements in $B$. This means $g$ \underline{is not} a function. However, on the right, we have $f:A\to B$ that assigns each element in $A$ only one element in $B$. This means $f$ \underline{is} a function.
\begin{center}
    \resizebox{.3\textwidth}{!}{%
\begin{circuitikz}
\tikzstyle{every node}=[font=\Large]
\draw  (3,11) ellipse (1cm and 3.25cm);
\draw  (7.25,11) ellipse (1cm and 3.25cm);
\draw  (3,13) circle (0.25cm);
\draw  (3,10) circle (0.25cm);
\node [font=\LARGE] at (3.25,12.5) {};
\node [font=\LARGE] at (3.25,12.5) {};
\draw  (2.75,10.25) circle (0cm);
\draw  (2.75,10) circle (0cm);
\draw  (2.75,10) circle (0cm);
\draw  (3,8.75) circle (0.25cm);
\draw  (3,11.5) circle (0.25cm);
\draw  (7,13) rectangle (7.5,12.5);
\draw  (7,11.75) rectangle (7,11.75);
\draw  (7,11.75) rectangle (7.5,11.25);
\draw  (7.25,10.5) rectangle (7.25,10.5);
\draw  (7,10.25) rectangle (7.5,9.75);
\draw  (7.25,9.25) rectangle (7.25,9.25);
\draw  (7,9) rectangle (7.5,8.5);
\draw [->, >=Stealth] (3.25,13) -- (7,12.75);
\draw [->, >=Stealth] (3.25,11.5) -- (7,11.5);
\draw [->, >=Stealth] (3.25,11.5) -- (7,10);
\draw [->, >=Stealth] (3.25,10) -- (7,8.75);
\draw [->, >=Stealth] (3.25,8.75) -- (7,8.75);
\node [font=\Large] at (5.25,15) {$g:A\to B$};
\node [font=\Large] at (3,7.25) {A};
\node [font=\Large] at (7.25,7.25) {B};
\end{circuitikz}
}%
~~~~
\resizebox{.3\textwidth}{!}{%
\begin{circuitikz}
\tikzstyle{every node}=[font=\Large]
\draw  (3,11) ellipse (1cm and 3.25cm);
\draw  (7.25,11) ellipse (1cm and 3.25cm);
\draw  (3,13) circle (0.25cm);
\draw  (3,10) circle (0.25cm);
\node [font=\LARGE] at (3.25,12.5) {};
\node [font=\LARGE] at (3.25,12.5) {};
\draw  (2.75,10.25) circle (0cm);
\draw  (2.75,10) circle (0cm);
\draw  (2.75,10) circle (0cm);
\draw  (3,8.75) circle (0.25cm);
\draw  (3,11.5) circle (0.25cm);
\draw  (7,13) rectangle (7.5,12.5);
\draw  (7,11.75) rectangle (7,11.75);
\draw  (7,11.75) rectangle (7.5,11.25);
\draw  (7.25,10.5) rectangle (7.25,10.5);
\draw  (7,10.25) rectangle (7.5,9.75);
\draw  (7.25,9.25) rectangle (7.25,9.25);
\draw  (7,9) rectangle (7.5,8.5);
\draw [->, >=Stealth] (3.25,13) -- (7,12.75);
\draw [->, >=Stealth] (3.25,11.5) -- (7,11.5);
\draw [->, >=Stealth] (3.25,10) -- (7,8.75);
\draw [->, >=Stealth] (3.25,8.75) -- (7,8.75);
\node [font=\Large] at (5.25,15) {$f:A\to B$};
\node [font=\Large] at (3,7.25) {A};
\node [font=\Large] at (7.25,7.25) {B};
\end{circuitikz}
}%
\end{center}

\end{example}

\begin{definition}[Injection, Surjection, and Bijection]\hfill

        Firstly, a function $f$ is \emph{injective} if distinct elements in $A$ are mapped to distinct elements in $B$. Secondly, $f$ is \emph{surjective} if every element in $B$ has at least one element in $A$ that gets mapped to it. Lastly, if $f$ is both injective and surjective, then $f$ is \emph{bijective}.
\end{definition}
\begin{center}
    \resizebox{.7\textwidth}{!}{%
\begin{circuitikz}
\tikzstyle{every node}=[font=\LARGE]
\draw  (3,11) ellipse (1cm and 3.25cm);
\draw  (7.25,11) ellipse (1cm and 3.25cm);
\draw  (3,13) circle (0.25cm);
\draw  (3,10) circle (0.25cm);
\node [font=\LARGE] at (3.25,12.5) {};
\node [font=\LARGE] at (3.25,12.5) {};
\draw  (2.75,10.25) circle (0cm);
\draw  (2.75,10) circle (0cm);
\draw  (2.75,10) circle (0cm);
\draw  (3,8.75) circle (0.25cm);
\draw  (3,11.5) circle (0.25cm);
\draw  (7,12.5) rectangle (7.5,12);
\draw  (7,11.75) rectangle (7,11.75);
\draw  (7,11.5) rectangle (7.5,11);
\draw  (7.25,10.5) rectangle (7.25,10.5);
\draw  (7,10.25) rectangle (7.5,9.75);
\draw  (7.25,9.25) rectangle (7.25,9.25);
\draw  (7,9) rectangle (7.5,8.5);
\draw [->, >=Stealth] (3.25,13) -- (7,13.5);
\draw [->, >=Stealth] (3.25,11.5) -- (7,11.25);
\draw [->, >=Stealth] (3.25,10) -- (7,10);
\draw [->, >=Stealth] (3.25,8.75) -- (7,8.75);
\node [font=\Large] at (3,7.25) {A};
\node [font=\Large] at (7.25,7.25) {B};
\node [font=\Large] at (5.25,15.25) {injective};
\draw  (7,13.75) rectangle (7.5,13.25);
\end{circuitikz}
~~
\begin{circuitikz}
\tikzstyle{every node}=[font=\Large]
\draw  (3,11) ellipse (1cm and 3.25cm);
\draw  (7.25,11) ellipse (1cm and 3.25cm);
\draw  (3,13) circle (0.25cm);
\draw  (3,10) circle (0.25cm);
\node [font=\LARGE] at (3.25,12.5) {};
\node [font=\LARGE] at (3.25,12.5) {};
\draw  (2.75,10.25) circle (0cm);
\draw  (2.75,10) circle (0cm);
\draw  (2.75,10) circle (0cm);
\draw  (3,8.75) circle (0.25cm);
\draw  (3,11.5) circle (0.25cm);
\draw  (7,11.75) rectangle (7,11.75);
\draw  (7.25,10.5) rectangle (7.25,10.5);
\draw  (7,11.25) rectangle (7.5,10.75);
\draw  (7.25,9.25) rectangle (7.25,9.25);
\draw  (7,9) rectangle (7.5,8.5);
\draw [->, >=Stealth] (3.25,13) -- (7,12.75);
\draw [->, >=Stealth] (3.25,11.5) -- (7,11);
\draw [->, >=Stealth] (3.25,10) -- (7,11);
\draw [->, >=Stealth] (3.25,8.75) -- (7,8.75);
\node [font=\Large] at (3,7.25) {A};
\node [font=\Large] at (7.25,7.25) {B};
\node [font=\Large] at (5.25,15.25) {surjective};
\draw  (7,13) rectangle (7.5,12.5);
\end{circuitikz}
~~
\begin{circuitikz}
\tikzstyle{every node}=[font=\Large]
\draw  (3,11) ellipse (1cm and 3.25cm);
\draw  (7.25,11) ellipse (1cm and 3.25cm);
\draw  (3,13) circle (0.25cm);
\draw  (3,10) circle (0.25cm);
\node [font=\LARGE] at (3.25,12.5) {};
\node [font=\LARGE] at (3.25,12.5) {};
\draw  (2.75,10.25) circle (0cm);
\draw  (2.75,10) circle (0cm);
\draw  (2.75,10) circle (0cm);
\draw  (3,8.75) circle (0.25cm);
\draw  (3,11.5) circle (0.25cm);
\draw  (7,11.75) rectangle (7,11.75);
\draw  (7.25,10.5) rectangle (7.25,10.5);
\draw  (7,11.75) rectangle (7.5,11.25);
\draw  (7.25,9.25) rectangle (7.25,9.25);
\draw  (7,9) rectangle (7.5,8.5);
\draw [->, >=Stealth] (3.25,13) -- (7,12.75);
\draw [->, >=Stealth] (3.25,11.5) -- (7,11.5);
\draw [->, >=Stealth] (3.25,10) -- (7,10);
\draw [->, >=Stealth] (3.25,8.75) -- (7,8.75);
\node [font=\Large] at (3,7.25) {A};
\node [font=\Large] at (7.25,7.25) {B};
\node [font=\Large] at (5.25,15.25) {bijective};
\draw  (7,13) rectangle (7.5,12.5);
\draw  (7,10.25) rectangle (7.5,9.75);
\end{circuitikz}
}%
\end{center}
For combinatorial purposes, one nice way to think of a function $f:A\to B$ as the placement of a set of objects $A$ into a set of boxes $B$:
    \begin{itemize}
        \item If no two objects are placed in the same box, $f$ is injective. 
        \item If no box is left empty, $f$ is surjective.
        \item If both occur, $f$ is bijective.
    \end{itemize}
We can use this thinking to phrase the following statements:\\
Let $f:A\to B$ be a function that assigns a set of objects $A$ to be placed into a set of boxes $B$.
\begin{itemize}
    \item $f$ is injective means $f$ ensures no two objects are placed in the same box.
    \item $f$ is surjective means $f$ leaves no box empty.
    \item $f$ is bijective means $f$ leaves no box empty and ensures no two objects are placed in the same box. 
\end{itemize}

But, why do we care about this as combinatorialists?
\begin{thm}\hfill

Again, let $f:A\to B$ as defined above.
    \begin{enumerate}\label{injsurjbij}
        \item If $f$ is injective, $|A|\leq |B|$.
        \item If $f$ is surjective, $|A|\geq |B|$.
        \item If $f$ is bijective, $|A|=|B|$
    \end{enumerate}
    where $|\cdot|$ denotes the number of elements in a set.
\end{thm}
First, an outline of a proof.
\begin{proof}\hfill

Let $f$ be a function such that $f:A\to B$ for sets $A$ and $B$.\\
    \begin{enumerate}
        \item If $f$ is injective, it ensures no two objects are in the same box. So, there must be at least enough boxes to house these distinctly assigned objects ($|A|\leq |B|$).
        \item If $f$ is surjective, it leaves no box empty. So, there must be at least enough objects to fill each box ($|A|\geq |B|$).
        \item If $f$ is bijective, it ensures no two objects are in the same box \emph{and} no box is left empty. As such, we know $|A|\leq |B|$ \emph{and} $|A|\geq |B|$ (i.e. that $|A|=|B|$).
    \end{enumerate}
\end{proof}
This simple statement is incredibly powerful. What this theorem firstly tells us is that there is a way to infer information about the \emph{sizes} of the sets when we have a specific property about the \emph{relationship} between those sets. Most importantly, the theorem also states that if there exists a bijection between two sets, they must be the exact same size. We can exploit this fact to enumerate easily! In fact, we already have.\\
Think back to our second problem: “how many integers are there between 28 and 352 inclusive?”. Perhaps we write out ``28 is $1^{st}$, 29 is $2^{nd}$,$\dots$'' and so on until we reach 352 or count on our 325 fingers one by one. Either way we know this would be annoying and tedious to count explicitly. By defining a bijective function between sets, we eased the trouble, $f(a)=a-27$ for $a\in \{28,29,30,\dots, 325\}$.\\
From this final ``proof'' in \ref{injsurjbij}, we can infer a lot of information. For clarity, if $f$ is bijective, each object has a distinct box \emph{and} each box has a distinct object. If we wish to reverse the process of assignment $f$ has committed, we could easily do so: simply reverse each assignment! \\
This thinking leads to the following theorem:
\begin{thm}\hfill

    The function $f:A\to B$ is a bijection if and only if $f^{-1}$ exists.
\end{thm}
This allows us to formalize our results:
\begin{itemize}
    \item Let $A=\{28,29,\dots, 352\}$, $B=\{1,2,\dots, 325\}$, and $f:A\to B$ be defined as $f(a)=a-27=b$ for $a\in A$ and $b\in B$.
    \item Since $f$ subtracts 27, $f^{-1}$ might add 27 ($f(a)=a-27=b\implies a=f(a)+27\implies f^{-1}(b)=b+27$).
    \item Thus, $f$ must be a bijection since there exists an $f^{-1}$. So, we know $|A|=|B|$. 
    \item Therefore, $|\{1,2,\dots, 325\}|=|\{28,29,\dots, 352\}|=325$.
\end{itemize}
We can use the same process with problems 3 and 4 of \ref{question1}
:
\begin{itemize}
\item Let $A=\{28,30, 32,\dots, 352\}$ and define $f(a)=\frac{1}{2}a-13=b$ for $a\in A$. What is $f^{-1}$? What is $B$? Conclude $|A|=|B|$.
\item Let $A=\{30,33,36,\dots,348,351\}$ and define $f(a)=\frac{1}{3}a=b$ for $a\in A$. What is $f^{-1}$? What is $B$? Conclude $|A|=|B|$.
\item (Extra!) Let $A=\{m,m+1,m+2,\dots, n-1, n\}$ and define $f(a)=\frac{1}{k}a-\ell$ for $m<k<n$ and $\ell\geq 0$. What is $f^{-1}$? What is $B$? Conclude $|A|=|B|$.
\end{itemize}

\section{Summary}
What we've done so far is introduce important techniques we often use in combinatorics:
\begin{itemize}
        \item  Let's simplify our current problem to a problem we already know how to solve
        \item  Now that we have found a method to solve a problem, let us generalize  
    \end{itemize}
These often guide thought in mathematical problem solving and we will utilize both in combinatorics!

\chapter{Introduction to counting techniques}

\subsection{Warmup}
Consider the following problem:
\begin{question}\label{question2}\hfill

    Physics 533 and Math 325 occur at the same time (12:10-1PM MWF). Phys 533 has 5 people enrolled while Math 325 has 21. 
    \begin{enumerate}
        \item How many people are in Math 325 or Physics 533 from 12:10-1 MWF? 
        \item How many people are in Math 325 and Physics 533 from 12:10-1 MWF? 
    \end{enumerate}
\end{question}
This may seem immediately obvious to some, but this highlights key counting principles used in combinatorics. 
\begin{answer}\hfill
    \begin{enumerate}
        \item There are 21 people in Math 325. There are 5 in Phys 533. So, there are 21+5 in either during the class period.
        \item Nobody can be in two places at once, so there must be zero people in both classrooms during the class period.
    \end{enumerate}
\end{answer}
This is a nice set up. We know the set of students in 325 is completely separate from the set of students in 533. If we use set theoretic language, the set of students in 325 is \emph{disjoint} from the set of students in 533. More on this later.\\
For now, consider yet another problem:
\begin{question}\hfill

    How many ways are there to flip a coin then a six sided dice? More generally, rolling an $m$ sided dice then an $n$ sided dice?
\end{question}
You might note that my first roll \emph{does not affect} my second roll. They are independent! Intuitively, if I roll heads, I have the opportunity to roll 1-6 on the dice. Likewise if roll tails. As such:
\begin{answer}\hfill

    There are 2 options for my first roll. Since each roll is independent, there are 6 options for my second roll. So, there must be $2\cdot 6=12$ ways to do this. More generally, since each roll is independent, there are $m\cdot n$ ways to do this with an $m$ sided dice and an $n$ sided dice.
\end{answer}

\section{More Set Theory}
\begin{definition}[Subset]\hfill

    Let $B$ be some set. Another set $A$ is a \emph{subset} of $B$ if every element of $A$ is an element of $B$ (for every $a\in A$, $a\in B$). We denote this relationship by $A\subseteq B$ if $A$ could equal $B$ itself. If $A$ cannot equal $B$, we write $A\subset B$.  
\end{definition}
\begin{definition}[Set Complement]\hfill

    Let $A\subseteq S$. The \emph{complement} of $A$ with respect to $S$, $A^C=S\setminus A$, is the set of elements in $S$ that are not in $A$ ($A^C=\{s\in S|s\notin A\}$). The complement of $S$ with respect to itself, $S^C=S\setminus S$, is the set with zero elements, $\emptyset=\{\}$. 
\end{definition}
  We can think of the set complement $S\setminus A$ as considering all the elements of $S$ that are not in $A$. If there are no such elements (i.e. $A\nsubseteq S$), $S\setminus A$ is $\emptyset$.\\
\begin{example}\hfill

Let $A=\{$integers between 28 and 352 divisible by 3$\}$ and $S=\{$integers between 28 and 352 inclusive$\}$
    \begin{itemize}
        \item $A\subset S$
        \item $A^C=S\setminus A=\{$integers between 28 and 352 not divisible by 3\}
    \end{itemize}
    
\end{example}
  
    \begin{rem}\hfill
    
        The choice of $S$ is important. If $S=\mathbb{Z}$ (the set of all integers), then $A^C$ in the previous example would become the set of \emph{all} integers excluding only integers between 28 and 352 that are divisible by 3. 
    \end{rem}

\begin{question}\hfill

    Let $A\subseteq S$. Is the function $f(A)=A^C$ a bijection? If so, what is $f^{-1}$?
\end{question}
\begin{answer}\hfill

$f$ assigns subset $A$ to its complement $A^C=S\setminus A$. 
    \begin{itemize}
        \item Let us say $f(A)=A^C=B$.
        \item By the definition of a set's complement, $B$ contains every element of $S$ that is not in $A$. Likewise, $A$ contains every element of $S$ that is not in $B$.
        \item In other words, $A^C=B$ \emph{and} $A=B^C$.
        \item So, $f^{-1}(B)=B^C$. Hence, $f^{-1}$ exists and $f$ is a bijection!
    \end{itemize}
\end{answer}
\vspace{.5em}
This is a nice property. Consider the problem of enumerating $k$-element subsets of $S=\{1,2,3,4,\dots, n\}$ ($k<n$).
    \begin{itemize}
        \item Let $A$ be the set of all $k$-elements subsets $A_i\subset S$. Moreover, let $B$ be the set of all the complements of elements in $A$.
        \item Then for each $A_i$, there is some set $B_i$ such that $f(A_i)=A_i^C=B_i\subset S$ of size $n-k$. 
        \item Since $f(A_i)=A_i^C$ is a bijection, the set of all $k$-element subsets, $A$, is the same size as the set of all ($n-k$)-element subsets, $B$. That is, $|A|=|B|$.
    \end{itemize}
    We will talk about what the actual sizes $|A|$ and $|B|$ later.

\begin{definition}[Set Union]\hfill

    The \emph{union} of two sets $A$ and $B$, denoted $A\cup B$, is the set of all elements in $A$ or $B$ ($A\cup B=\{x|x\in A$ or $x\in B$\}).
\end{definition}
\begin{definition}[Set Intersection]\hfill

    The \emph{intersection} of two sets $A$ and $B$, denoted $A\cap B$, is the set of all elements in $A$ and $B$ ($A\cap B=\{x|x\in A$ and $x\in B$\}).
\end{definition}
\begin{definition}[Disjointness]\hfill

    If the intersection between sets $A$ and $B$ is empty, denoted $A\cap B=\emptyset$, we say $A$ and $B$ are \emph{disjoint}.
\end{definition}
\begin{example}\hfill

Let $A=\{$the set of students in Math 325$\}$ and $B=\{$the set of students in Physics 533$\}$.
\begin{itemize}
    \item $A\cup B=\{$the set of students in either couse$\}$
    \item $A\cap B=\emptyset$ since now student is in both classes
\end{itemize}
\end{example}

\section{Addition \& Multiplication Principle}
Thinking back to the initial problem of this section (\ref{question2}), we could simply add the number of people in 325 and 533 to get the total number \emph{since no student is in both classes}. This principle of complete separation allowing simple addition can be formalized: 
\begin{thm}[Addition Principle]\hfill

    If $A\cap B=\emptyset$, $|A\cup B|=|A|+|B|$.\\
    More generally, if $A_i\cap A_j=\emptyset$ for all $i\neq j$ and $1\leq i,j\leq n$, then $|A_1\cup A_2\cup \cdots \cup A_n|=|A_1|+|A_2|+\cdots +|A_n|$.
\end{thm}


\begin{definition}[Cartesian Product]\hfill

    Let $A\times B$ be the set of ordered pairs $(a,b)$ where $a\in A$ and $b\in B$.
\end{definition}
\begin{thm}[Multiplication Principle]\hfill

    Let $S=A\times B$. Then, $|S|=|A\times B|=|A|\cdot|B|$.\\ 
    More generally, let $S=A_1\times A_2\times \cdots \times A_n$. Then, $|S|=|A_1\times A_2\times \cdots \times A_n|=|A_1||A_2|\cdots|A_n|$
\end{thm}
\begin{example}\hfill

    From warmup problem 2, let $A=\{H,T\}$ be the results from flipping a coin and $B=\{1,2,3,4,5,6\}$ be the results from rolling the six-sided dice. Then,
    \begin{center}
    $A\times B=\{(H,1),(H,2),(H,3),\dots,(T,5),(T,6)\}$
    \end{center}
    and $|S|=|A\times B|=|A||B|=2\cdot 6=12$.
\end{example}


\subsection{(In)dependence}
Note that the \emph{independence} of warmup problem 2 is important. Flipping a coin \emph{has no impact} on the rolling of the dice. Consider a different problem:
\begin{question}\hfill

    There are four orbs labelled $A,B,C,D$ in an urn. You will draw from the urn four times. For each draw, you take a single orb out, write down its label, and replace it in the urn. How many possible label sequences could you write down?
\end{question}
\begin{answer}
    
    \begin{itemize}
        \item Let $A_1, A_2, A_3,$ and $A_4$ be the sets of possible draws on draw 1, 2, 3, and 4 respectively and $S$ be the set of possible label sequences you could write down.  
        \item Since all four balls in the urn before you draw first, $A_1=\{A,B,C,D\}$.
        \item Moreover, since you placed the orb you drew back into the urn, you have a chance to draw any of the four orbs again: $A_2=\{A,B,C,D\}$.
        \item Likewise with $A_3$ and $A_4$.
        \item Thus, the total possible label sequences there are is: $|S|=|A_1\times A_2\times A_3\times A_4|=|A_1||A_2||A_3||A_4|=4^4=256$.
\end{itemize}
\end{answer}


    Consider a small change to your process:
\begin{question}\hfill

        You really like writing down labels, so you continue with your setup from the previous problem. However, this time instead of replacing the orb after drawing it, you set it to the side. You then draw another orb and set it next to the previously drawn orb. You repeat this process until there are no orbs left. Then, you write down the sequence of labels you see. How many possible label sequences could you see?
    \end{question}

\begin{answer}\hfill

    \begin{itemize}
        \item Let $B_1, B_2, B_3,$ and $B_4$ be the sets of possible draws on draw 1, 2, 3, and 4 respectively and $T$ be the set of possible label sequences you could write down. 
        \item Since all four balls in the urn before you draw first, $|B_1|=4$.
        \item Whatever you drew, it stays outside of the urn. Since there are 3 orbs left, you have 3 options for your next draw: $|B_2|=3$.
        \item Continuing on, you have taken yet another option for your third draw away: $|B_3|=2$.
        \item Lastly, since you drew one of the two remaining orbs, that leaves one left for your final draw $|B_4|=1$.
        \item Thus, $|T|=|B_1\times B_2\times B_3\times B_4|=|B_1||B_2||B_3||B_4|=4\cdot 3\cdot 2\cdot 1=24$.
    \end{itemize}
\end{answer}
    So, replacing the orbs every time gave us a count of 256 while leaving them out every time gave us a count of 24. What happened? One way to think of this is that by replacing the orbs, you have allowed for \emph{repetitions} to exist in your set $S$ (e.g. by drawing $A$ then $A$ then $C$ then $A$). Contrarily, by leaving the orbs out after drawing, you remove the ability to have repeats. As such, you remove every element in $S$ that has any repeat labels to obtain a new set $T$, the set of label sequences with \emph{distinct} labels.
\begin{question}[To Ponder]\hfill

        Let $S$ and $T$ be defined as in the above problem. What is $T^C=S\setminus T$? How big is $T^C$ in this case?
    \end{question}
    \begin{question}[To Ponder]\hfill
    
        Let $S$ be defined as in the above problem. How many label sequences have only 2 repeats? i.e. sequences like $(A,A,B,C)$ or $(D,A,C,D)$.
    \end{question}
    \begin{question}[To Ponder]\hfill
    
        Consider drawing $n$ labelled orbs from an urn with replacement. How many ways can you do this? Without replacement?
    \end{question}

\section{Summary}
What we've covered here:
    \begin{itemize}
        \item A function $f:A\to B$ having an inverse function $f^{-1}:B\to A$ means $|A|=|B|$.
        \item Set complements, unions, and intersections. 
        \item Addition principle: the size of a union of disjoint sets is the sum of their sizes.
        \item Multiplication principle: the size of a product of two sets is the product of their sizes.
        \item The difference between counting with dependence versus without.
    \end{itemize}
\end{document}